% SOURCE_AUTHORITY: sga5_fr_workpass.tex, lines 14118--14143 (LF-normalized unit).
% SOURCE_SHA256: 791F4EFFC5E02832D5D77ED03518C8156D6F07E4C8238B03545DB93D883FBB28
\section*{7. Comentários sobre as condições de validade da fórmula de Lefschetz}

\subsection*{7.1}
Naturalmente, na prática, parte-se de um esquema $X$ próprio sobre $k$ e de um endomorfismo
\[
(u, f) : (X, F) \longrightarrow (X, F) \quad ,
\]
tal que:
\begin{itemize}
\item $F$ é um feixe construtível de $A$-módulos,
\item o conjunto $X^f$ dos pontos fixos de $f$ é finito.
\end{itemize}

O problema consiste em definir, sob estas condições, termos locais $\tau_x(u, f)$, $x \in X^f$, tais que
\[
\text{(*)} \qquad \Tr_A(\R\Gamma_X(u, f)) = \sum_{x \in X^f} \tau_x(u, f) \quad .
\]

Usando os dois teoremas de dualidade, global e local, Verdier consegue definir os termos $\tau_x(u, f)$ de modo que a fórmula (*) seja verdadeira (cf. III). Infelizmente, a determinação explícita dos $\tau_x(u, f)$ definidos por Verdier em termos de invariantes locais conhecidos parece difícil, daí o interesse das definições propostas no capítulo anterior.\footnote{Veja-se, contudo, III B para a comparação dos termos locais de 6.3.1 com os termos locais de Verdier.}

Em contrapartida, estas últimas definições supõem hipóteses suplementares sobre $F$, $f$ e $u$, a saber, a existência de $U, U', G, f', \varphi$ com as propriedades exigidas em 5.1, 6.1, 6.2.
